\begin{pastexam}[label={px:nccu-106-dsa-02}]{政大資科所 106 學年度 第 2 題}
（問答）

(10\%) Consider the following function \texttt{f}. How many lines, as a function of \texttt{n}
(in $\Theta$(\texttt{n}) form), does \texttt{f(n)} print? Write a recurrence and solve it. You may
assume \texttt{n} is a power of 2.
\begin{minted}[fontsize=\small]{c}
void f(int n){
    if (n > 1){
        printf("still going\n");
        f(n / 2);
        f(n / 2);
    }
}
\end{minted}
\end{pastexam}

\begin{pastsoln}
答案：$f(n)$ 印出 $n - 1$ 行，是 $\Theta(n)$。

題目的「in $\Theta(\texttt{n})$ form」是要用 $\Theta$ 記號寫出答案的量級。

\textbf{列遞迴式}（資結書第~2.1~節）。 令 $P(n)$ 是 \texttt{f(n)} 印出的行數。 $n = 1$ 時
\texttt{if} 不成立，什麼都不印；$n > 1$ 時自己印一行，再呼叫兩次 \texttt{f(n/2)}，各印
$P(n/2)$ 行。 $n$ 是 $2$ 的冪次，$n/2$ 沒有捨去的問題。 所以
\[
P(1) = 0, \qquad P(n) = 2P(n/2) + 1 \quad (n \ge 2).
\]

\textbf{解遞迴式}。 令 $n = 2^m$，一直展開：
\begin{align*}
P(n) &= 2P(n/2) + 1 \\
     &= 2\bigl(2P(n/4) + 1\bigr) + 1 = 4P(n/4) + 2 + 1 \\
     &= 8P(n/8) + 4 + 2 + 1 \\
     &\;\;\vdots \\
     &= 2^k P(n/2^k) + (2^{k-1} + \cdots + 2 + 1) = 2^k P(n/2^k) + 2^k - 1.
\end{align*}
取 $k = m = \lg n$，$P(n/2^m) = P(1) = 0$，得
\[
P(n) = 0 + 2^m - 1 = n - 1.
\]
換成遞迴樹的說法：第 $j$ 層（$j = 0, 1, \ldots, m - 1$）有 $2^j$ 個呼叫，參數都大於 $1$，各印一行；
第 $m$ 層的 $2^m$ 個呼叫參數是 $1$，不印。 總共 $1 + 2 + \cdots + 2^{m-1} = 2^m - 1$ 行。

\textbf{用數學歸納法確認} $P(n) = n - 1$。 $n = 1$：$P(1) = 0 = 1 - 1$。 設 $P(n/2) = n/2 - 1$，則
$P(n) = 2(n/2 - 1) + 1 = n - 1$。 例如 $n = 1, 2, 4, 8$ 分別印 $0, 1, 3, 7$ 行。

所以 $P(n) = n - 1 = \Theta(n)$（$n \ge 2$ 時 $n/2 \le n - 1 \le n$）。

也可以直接用主定理（資結書定理~2.4）：$a = 2$、$b = 2$，$n^{\log_b a} = n$；$f(n) = 1 =
O(n^{1 - \epsilon})$，取 $\epsilon = 1$，是第一種情形，$P(n) = \Theta(n)$。 主定理只給量級，題目
要求「solve it」，寫出 $P(n) = n - 1$ 比較完整。
\end{pastsoln}
